\section{Limitations and honest negatives}
\label{sec:limitations}

We separate what the study establishes from what it does not.

\subsection{Cohort scope and the cross-cohort caveat}
\label{sec:lim:caveat}

The head-to-head in Table~\ref{tab:krasnov} is a \emph{cross-cohort} estimate.
Our funnel keeps $591$ cell lines against the paper's $23$; we adopt their split
scheme and dedup rule exactly, but the underlying population differs. We can now
attribute the $0.0383$ gap: running the published model class on our identical
cohort and fold partition closes it to $-0.0040$, so the gap is cohort scope, not
model class. The claim we defend is therefore: under the published protocol, on
a broader and more heterogeneous cohort, a D-MPNN and the published LightGBM are
statistically indistinguishable, and both sit $0.038$ below the number the paper
reports on its narrower cohort.

\subsection{Small per-metal cohorts make part of the result unresolvable}
\label{sec:lim:small}

Three of the five metals have too few unique ligands to resolve an effect:
Rh $51$, Re $74$, Os $36$. Os is the clearest case: its within-metal AUC swings
from $0.3133$ to $0.5867$ across five seeds (spread $0.2734$), larger than any
method change would move it. That swing \emph{is} the unresolvability. We
therefore report Os, Rh and Re as measured-but-not-gated and draw no metal-level
conclusion from them individually. In particular, the sub-$0.5$ cross-metal
point estimates for Os ($0.4950$) and Re ($0.4429$) are \emph{not} a
below-chance effect: a label-permutation null on the fixed cross-metal ranking
places both inside the null (null SD $\pm0.10$ at $n=36$, $\pm0.09$ at $n=74$),
and only Ru is significantly above chance under the fingerprint scorer after
Holm correction (the message-passing backbone places Re above chance at $0.7368$,
Holm $p=0.010$; see
\S\ref{sec:results:collapse}). Re's three-seed cross-metal range
($0.4429$--$0.4798$) is tight, but a tight range around an unresolvable level is
still unresolvable.

\subsection{The pooled cross-metal mean is Ru-dominated}
\label{sec:lim:pooled}

Ru is $64\%$ of the deduplicated cohort. Any pooled cross-metal mean is largely
a Ru number, and quoting it without the per-metal table would overstate the
breadth of the finding. We quote per-metal throughout and treat the pooled value
as descriptive only.

\subsection{A cohort-construction trap in a related protocol}
\label{sec:lim:cohort}

A preprocessing filter applied to a pre-filtered artifact can manufacture a
below-chance result, so we record one we hit while building a related
leave-one-metal-out protocol. The default cohort file is a snapshot already
filtered to short exposure times ($\leq24$\,h); raising the threshold at training
time cannot recover the dropped rows, and the loaded cohort stayed
byte-identical --- the filter silently did nothing. The consequence is a
below-chance fold: the Rh held-out set collapses to $108$ test ligands and scores
$0.3812$, dragging the pooled LOMO mean to $0.5575$, with no change to the model.
Readers building on MetalCytoToxDB can hit the same trap. \MEASURED:
\texttt{metrics/by\_metric/metal\_lomo\_auc\_wasserstein.json}.

\subsection{Two MoleculeNet task gaps remain}
\label{sec:lim:bace}

Under a protocol-matched comparison (our own chemprop~v2 re-run, same scaffold
split and seeds) the D-MPNN wins bbbp ($+0.0133$) and hiv ($+0.0082$) and loses
bace ($0.7933 \pm 0.0250$ against $0.8458 \pm 0.0113$, $-0.0525$) and clintox
($0.8303$ against $0.8847$, $-0.0544$). Both losses exceed seed noise, so we
report them as honest negatives. They are generic ADMET endpoints outside this
paper's target chemistry and do not bear on the metal-axis claims; the backbone
is otherwise competitive (mean $0.8280$), which is what lets us argue in
\S\ref{sec:results:moleculenet} that the metal collapse is a property of the
axis and not of a broken model. Seed-ensembling the five runs is a free
inference-time lever that lifts the mean by $+0.0123$ on a pinned-split
replication (\texttt{l2\_moleculenet\_ensemble.json}) without closing the bace
gap.

\subsection{Leakage controls}
\label{sec:lim:leak}

Cross-metal evaluation is only meaningful if test ligands are absent from
training. The LOMO protocol drops any training ligand that appears in the
held-out metal, so the cross-metal column is ligand-disjoint by construction.
The within-metal column uses 5-fold CV over a metal's own ligands and shares no
ligands across folds by construction. The scaffold control in
Table~\ref{tab:transfer} holds out whole Bemis--Murcko scaffolds within a metal,
so its ``net metal'' column is the collapse with the ordinary novel-scaffold
penalty already removed. Two caveats remain: the scaffold-split control is a
single seed ($42$), and both the cross-metal and scaffold splits are
ligand-/scaffold-disjoint but not cell-line-disjoint.

\subsection{What we do not claim}
\label{sec:lim:notclaim}

\begin{itemize}
\item We do not claim the collapse is caused by any specific chemical mechanism.
We measure that it happens; we do not identify why.
\item We do not claim any metal's cross-metal AUC is \emph{below} chance under
the fingerprint scorer. Under its label-permutation null none is, and only Ru is
significantly above chance after Holm correction. The message-passing backbone
moves Re above chance ($0.7368$, Holm $p=0.010$); Ir's below-chance reading is
significant only as an uncorrected one-sided test ($p=0.038$, two-sided $0.076$,
Holm $0.23$) and does not survive the family correction, so we do not claim it.
The sub-$0.5$ fingerprint point estimates for Os and Re are unresolvable at
$n=36$--$74$.
\item We do not claim metal is the \emph{only} shift axis. A same-metal scaffold
split already costs $0.10$--$0.12$ on the two metals where the penalty resolves;
the metal axis is the largest single remaining shift, not the sole source of the
collapse.
\item We do not claim to beat the published model on its DOI-split replication:
we are $0.0383$ below it cross-cohort, and the same-cohort model-class gap is
$-0.0040$. We \emph{do} claim a within-metal win for Ir with the
fingerprint-plus-XGBoost scorer (Table~\ref{tab:transfer}: $+0.09$ to $+0.13$,
robust to both tie rules and every seed); Ru is at parity, decided by the tie
rule. The within-metal column is that scorer's, not the D-MPNN's: re-measured
on the same cohort and folds the D-MPNN scores $0.7041$ (Ru) and $0.6912$ (Ir)
against the GBDT's $0.8271$ and $0.8430$, i.e.\ below on all five metals.
\item We do not claim a tmQM pretraining benefit. On the fixed scaffold split,
measured on two metals, no Holm-corrected effect is significant: Ru's paired
delta is $-0.0094$ (t-based 95\% CI $[-0.0406,+0.0218]$), and Ir's is a
suggestive $+0.0486$ that does not clear the paired $t$-test ($3/5$ seeds
positive, CI includes zero, one-sided $p=0.068$; the 5-seed bootstrap's
$p=0.012$ is anti-conservative at that sample size). The earlier $+0.0181$ came
from a degenerate, molecule-level splitter and does not survive the fix.
\item We do not claim the bond-perception head wins. It loses macro-F1(type).
\item We do not claim Ru exceeds the published $0.81$ unqualified; that is
decided by the tie rule (\S\ref{sec:ablation:tie}).
\item We do not claim static metal conditioning mitigates the collapse. A metal
token, a metal-descriptor embedding and FiLM all leave the cross-metal mean at
or below the unconditioned baseline (\S\ref{sec:results:method}); the mitigation
that works is few-shot adaptation to the target metal, and it needs $20$--$40$
labelled ligands.
\end{itemize}

\section{Conclusion}
\label{sec:conclusion}

Metal identity is a first-class distribution shift for ligand-based activity
models. On a deduplicated MetalCytoToxDB cohort, a single shared scorer loses
$0.18$--$0.51$ ROC-AUC when it moves from a metal it has seen to one it has not,
and a same-metal scaffold-split control shows that only about half of that drop
is ordinary novel-scaffold generalisation. Under a Holm-corrected label-permutation null, and
at the cohort sizes available, the cross-metal signal is significantly above
chance only for Ru under the fingerprint scorer; an independent message-passing
backbone additionally recovers Re ($0.7368$, Holm $p=0.010$), which is why we report
Re as representation-dependent rather than as a failed transfer. The sub-$0.5$
fingerprint point estimates for Os and Re are not below-chance effects. tmQM
pre-training of the encoder does \emph{not} transfer to the activity head on a
genuine scaffold split, on either of the two metals tested (Ru $-0.0094$; Ir a
suggestive $+0.0486$ that does not clear the paired $t$-test, $p=0.068$), and
running the
published protocol's own model class on our
identical cohort and folds reproduces it to within $-0.0040$ --- so the honest
negative is about cohort scope, not model class.

The actionable reading is narrow but clear. A ligand-only QSAR model trained on
one metal should not be deployed to rank candidates for a different metal without
re-validation: the features that carry within-metal signal do not survive the
metal change, and feeding the model the metal identity does not repair this ---
static metal conditioning (token, descriptor, FiLM) leaves the cross-metal mean
at or below the unconditioned baseline. What does work is cheap: a handful of
labelled ligands from the target metal ($20$--$40$) recovers most of the collapse
($+0.058$ to $+0.091$ AUC, Holm $p<0.004$). We therefore recommend (i) few-shot
adaptation to the target metal rather than a metal-conditioned zero-shot model,
and (ii) reporting cross-metal leave-one-metal-out alongside the within-metal
number --- with a same-metal scaffold split as the control --- whenever a
dataset spans more than one metal. Cross-metal transfer deserves to be a
reported benchmark axis, not an implicit assumption.

\section*{Data and code availability}

All cohorts, configs and scripts are versioned in the accompanying repository.
Every quantitative claim in this paper cites a JSON artifact under
\texttt{metrics/by\_metric/}, listed at the corresponding table or figure.
